\chapter{Methodisches Vorgehen}

Zur Beantwortung der Forschungsfragen wird ein mehrstufiges Vorgehen verfolgt, das Literaturrecherche, konzeptionelle Modellierung sowie eine prototypische Umsetzung kombiniert.

\begin{itemize}


\item \textbf{Literaturrecherche:} 

Der Forschungsstand zu Datenqualität, Metadatenqualität, Linked Data Quality sowie bestehenden Bewertungsmodellen für RDF-basierte Metadaten wird systematisch analysiert. Zur Identifikation relevanter Arbeiten und zur Strukturierung der Literatur werden KI-gestützte Recherche- und Analysewerkzeuge eingesetzt. Diese unterstützen insbesondere bei der Exploration thematischer Zusammenhänge, der Identifikation relevanter Forschungsarbeiten sowie der systematischen Aufbereitung bestehender Bewertungsansätze.

\item \textbf{Entwicklung eines konzeptionellen Bewertungsmodells:}  
Auf Basis der Literatur und der Analyse bestehender SHACL-Validierungsergebnisse wird ein konzeptionelles Bewertungsmodell für DCAT-AP.de Metadaten entwickelt. Ein zentraler Bestandteil dieses Modells ist die Untersuchung, inwiefern eine agentenbasierte Architektur geeignet ist, regelbasierte Validierungssysteme um kontextuelle Analysefähigkeiten zu erweitern. Dabei wird betrachtet, wie spezialisierte KI-Agenten Validierungsergebnisse interpretieren, zusätzliche Kontextinformationen berücksichtigen und daraus verbesserte Qualitätsbewertungen ableiten können.

\item \textbf{Prototypische Umsetzung und Evaluation:}  
Das entwickelte Konzept wird prototypisch umgesetzt und auf repräsentative Metadatensätze angewendet, um die praktische Umsetzbarkeit sowie den Nutzen der vorgeschlagenen Architektur zu evaluieren.

\item \textbf{Evaluierung und Reflexion:} Abschließend werden das entwickelte Bewertungsverfahren sowie die Gewichtungsentscheidungen kritisch reflektiert, methodische Limitationen aufgezeigt und mögliche Erweiterungen diskutiert.
\end{itemize}

\chapter{Gliederung}

\begin{enumerate}
    \item Einleitung
    \begin{itemize}
        \item Motivation und Problemstellung
        \item Zielsetzung und Forschungsfragen
        \item Aufbau der Arbeit
    \end{itemize}

    \item Grundlagen
    \begin{itemize}
        \item Open Government Data und GovData
        \item DCAT-AP.de, RDF und Linked Data
        \item SHACL und Konzepte der Metadatenqualität
    \end{itemize}

    \item Stand der Forschung
    \begin{itemize}
        \item Daten- und Linked-Data-Qualitätsmodelle
        \item Bewertungsframeworks für Metadaten
        \item KI-gestützte Ansätze zur Metadatenanalyse
    \end{itemize}

    \item Konzept eines Metadaten-Qualitätsmodells
    \begin{itemize}
        \item Definition relevanter Qualitätsdimensionen
        \item Ableitung von Qualitätsmetriken
        \item Aggregation und Scoring der Validierungsergebnisse
    \end{itemize}

    \item Architektur und prototypische Umsetzung
    \begin{itemize}
        \item Verarbeitung von SHACL-Validierungsreports
        \item Systemarchitektur und Datenfluss
        \item Implementierung der Qualitätsmetriken und Reporting
    \end{itemize}

    \item Evaluation
    \begin{itemize}
        \item Datengrundlage und Versuchsaufbau
        \item Anwendung des Bewertungsmodells
        \item Analyse und Diskussion der Ergebnisse
    \end{itemize}

    \item Fazit und Ausblick
\end{enumerate}

\chapter{Vorläufiger Zeitplan}

\begin{center}
\resizebox{\textwidth}{!}{
\begin{ganttchart}[
    time slot format=isodate,
    compress calendar,
    x unit=0.07cm,
    y unit chart=0.8cm,
    vgrid={*1{draw=gray!40, dotted}},
    hgrid,
    bar height=0.40,
    group height=0.40,
    milestone height=0.40,
    bar/.append style={fill=teal!60, draw=teal!80!black},
    group/.append style={fill=black, draw=black},
    milestone/.append style={fill=black, draw=black},
    title/.append style={draw=black, fill=white},
    title height=1,
    title label font=\bfseries\footnotesize,
    bar label font=\footnotesize,
    bar label node/.append style={
    text width=5cm,
    align=left
    },
    group label node/.append style={
        text width=5cm,
        align=left
    },
    milestone label node/.append style={
        text width=2.5cm,
        align=center
    },
    group label font=\bfseries\footnotesize,
    milestone label font=\bfseries\footnotesize,
    calendar week text={\footnotesize\currentweek},
]{2026-04-01}{2026-09-30}

% Kopfzeile: Monate + Kalenderwochen
\gantttitlecalendar{month=name, week} \\

% --- Always-on Stränge ---
\ganttgroup{Durchgängig}{2026-04-01}{2026-09-30} \\
\ganttbar{Schreiben / Dokumentation (laufend)}{2026-04-01}{2026-09-30} \\
\ganttbar{Literaturrecherche (laufend, Schwerpunkt Apr--Mai)}{2026-04-01}{2026-05-31} \\

% --- Grundlagen / frühe Kapitel ---
\ganttgroup{Grundlagen \& Kontext}{2026-04-01}{2026-05-31} \\
\ganttbar{GovData/DCAT-AP.de Analyse}{2026-04-08}{2026-05-10} \\
\ganttbar{Stand der Forschung strukturieren}{2026-04-01}{2026-05-24} \\

% --- Iterative Zyklen (Konzept -> Implementierung -> Test/Eval) ---
\ganttgroup{Iteration 1 (v1)}{2026-06-01}{2026-06-28} \\
\ganttbar{Konzept v1: Dimensionen \& Metriken}{2026-06-01}{2026-06-07} \\
\ganttbar{Implementierung v1 (Basis-Prototyp)}{2026-06-08}{2026-06-18} \\
\ganttbar{Test/Eval v1 (kleine Stichprobe)}{2026-06-19}{2026-06-28} \\

\ganttgroup{Iteration 2 (v2)}{2026-06-29}{2026-07-26} \\
\ganttbar{Konzept v2: Scoring \& Gewichtung}{2026-06-29}{2026-07-05} \\
\ganttbar{Implementierung v2 (Metriken/Reports)}{2026-07-06}{2026-07-16} \\
\ganttbar{Test/Eval v2 (größere Stichprobe)}{2026-07-17}{2026-07-26} \\

\ganttgroup{Iteration 3 (v3)}{2026-07-27}{2026-08-23} \\
\ganttbar{Konzept v3: Verfeinerung \& Robustheit}{2026-07-27}{2026-08-02} \\
\ganttbar{Implementierung v3 (Stabilisierung)}{2026-08-03}{2026-08-13} \\
\ganttbar{Systematische Evaluation (repräsentativ)}{2026-08-14}{2026-08-23} \\

% --- Abschlussphase ---
\ganttgroup{Finalisierung}{2026-08-24}{2026-09-30} \\
\ganttbar{Ergebnisse aufbereiten (Tabellen/Plots/Reporting)}{2026-08-24}{2026-09-06} \\
\ganttbar{Diskussion \& Fazit finalisieren}{2026-09-07}{2026-09-20} \\
\ganttbar{Korrekturen, Review, Abgabevorbereitung}{2026-09-21}{2026-09-30} \\

% Meilensteine
\ganttmilestone{Review}{2026-09-20} \\
\ganttmilestone{Abgabe}{2026-09-30}

\end{ganttchart}
}
\end{center}